\documentclass[11pt,a4paper]{article}  % Déclare la classe du document.

\newcommand{\chaptermark}[2]{\og #2 \fg{} (#1)} % ligne uniquement pour éviter une erreur en cas d'article
\include{../1ere-preambule}

\usepackage{epic}   % extension et simplification de  l'environnement picture
\usepackage{graphicx}    % permet d'inclure des graphique externe dans un document
\usepackage{arydshln}   %permet d'avoir des lignes en pointillés

\newcommand{\lignes}[1]{
	\setlength{\unitlength}{1mm}
	\vskip 0,4cm\noindent
	\multido{\i=0+1}{#1}{%
		\noindent
		\begin{picture}(150,5)
			\linethickness{0,005mm}
			\put(-5,0){\dashline{1}(175,5)(5,5)}
		\end{picture}\vskip 0,001cm
	}
	\vskip 0cm
}

\rhead{DS n$^{\circ}03$ - Équations et inéquations}

\newcommand{\va}[1]{\lvert \ #1 \ \rvert}
\newcommand{\vaf}[1]{\left\lvert \ #1\  \right\rvert}

\begin{document}
	\graphicspath{{images/}}
	
	\noindent Nom: ${\scriptstyle  .........................................................}$ \hskip 1.8cm Prénom: ${\scriptstyle  ...................................................}$ 
	
	\vspace{.5cm}
	
	\begin{center}
		
		
		\fbox{
			\begin{minipage}[]{15cm}
				\begin{center}
					\LARGE{\rule{0.5cm}{0pt}\rule[-0.25cm]{0pt}{0.9cm}
						DS n$^{\circ}03$ $-$ Automatismes\\
						Équations et inéquations 
					}
				\end{center}
			\end{minipage}
		}
	\end{center}

\begin{center}\textbf{55min - Barème indicatif}\end{center}
\begin{center}\textbf{Les élèves avec un tier-temps ne traitent pas les questions avec  le symbole \ding{46}}\end{center}
\begin{center}\textbf{Calculatrice non autorisée.\\Les résultats devront être justifiés (à l'aide de calcul ou de propriétés).}\end{center}



\exerciceDS{ $-$ équation du premier degré $-$ (7 points)}

Résoudre dans $\mathbb{R}$ chacune des équations suivantes.
\begin{multicols}{3}
	\begin{enumMet}
		\item $5 x-3=7x+9$;
		\item \ding{46} $2 x-5=-4 x+1$;
		
		\columnbreak
		\item $2x+14=5x-4$;
		\item $0,3 x-0,7=0,5x+0,3$;
		
		\columnbreak
		\item $\dfrac{1}{3} x-\dfrac{2}{5}=-\dfrac{1}{4} x+\dfrac{5}{4}$;
	\end{enumMet}
\end{multicols}

\exerciceDS{ $-$ équation produit nul $-$ (4,5 points)}

Résoudre dans $\mathbb{R}$ chacune des équations suivantes.
\begin{multicols}{3}
	\begin{enumMet}
		\item $5(x-9)(x+4)=0$;
		\item $-3(3x-15)(-5x-20)=0$;
		\item \ding{46} $-15(7 x+7)(-x+4,5)=0$.
	\end{enumMet}
\end{multicols}

\exerciceDS{ $-$ Résoudre une inéquation du premier degré $-$ (6 points)}

Résoudre dans \R chacune des inéquations suivantes.
\begin{multicols}{3}
	\begin{enumMet}
		\item $-4 x +17 \geqslant 2x-7$;
		\item $-5 x+12 \leqslant 0$;
		\item $\dfrac{1}{3} x-\dfrac{1}{2} < \dfrac{1}{4}x-5$;
	\end{enumMet}
\end{multicols}

\exerciceDS{ $-$ équation de la forme $x^{2}=a$ $-$ (3 points)}

Résoudre dans \R les équations suivantes.
\begin{multicols}{3}
	\begin{enumMet}
		\item $x^{2}=9$;
		\item \ding{46} $x^{2}-8=0$;
		\item $3 x^{2}+6=0$;
	\end{enumMet}
\end{multicols}


\end{document}

